% ==============================================================================
% Frontmatter: Title Page, Preface, Methodology, and Study Guide
% ==============================================================================
\thispagestyle{empty}
\begin{center}
  \vspace*{2cm}
  {\fontsize{28}{34}\selectfont\bfseries\color{AcadNavy} Academic English Mastery}\\[1em]
  {\fontsize{16}{22}\selectfont\color{AcadSlate} A Comprehensive Textbook \& Interactive Workbook}\\[0.8em]
  {\fontsize{14}{20}\selectfont\color{AcadNavy}\bfseries 基于 \LaTeX\ 的学术英语教学书与实战练习册}\\[2em]
  {\color{AcadNavy}\rule{0.6\textwidth}{1.5pt}}\\[2.5cm]
  {\large\textbf{Academic English Curriculum Development Group}}\\[0.4em]
  {\normalsize\textbf{学术英语课程研制组}}\\[3cm]
  {\small\color{AcadGray} First Edition (2026) \quad\textbar\quad 第一版（2026年版）}\\[0.3em]
  {\footnotesize\color{AcadGray} Typeset with XeLaTeX under TeX Live 2026 System}
\end{center}
\clearpage

\chapter*{Preface \& Methodology / 前言与编写说明}
\addcontentsline{toc}{chapter}{Preface \& Methodology / 前言与编写说明}

本教材与配套实战练习册（\textit{Academic English Mastery}）严格对标国际主流高水平期刊（Nature, Science, IEEE/ACM Transactions, Cell, Lancet, AER）及顶级学术会议（NeurIPS, ICML, ICLR, CVPR, ACL）的真实发表级学术英语规范，旨在打破中国学者与高年级研究生在学术写作中长久存在的“语法正确但生硬冗长、词汇孤立但缺乏语境支配价、句式散乱且头重脚轻、段落堆砌而信息断裂、过度依赖 AI 生成空洞套话”五大核心痛点。

\section*{The Four Pedagogical Pillars / 四大学术语言工程支柱}
\begin{enumerate}[label=\textbf{\color{AcadNavy}\arabic*.}, leftmargin=2em]
  \item \textbf{词汇语境支配价与学术熟词生义（Lexical Valency \& Polysemy）}：不再孤立背诵单词，而是透彻解析 1,577 组高频学术多义词的核心学术义项、配价框架（Valency Frames）与强搭配网络（Collocations）。
  \item \textbf{句法骨架与抗头重脚轻工程（Syntactic Constituents \& End-Weight Principle）}：以五大核心句型为基础，贯彻“短主语引导、谓语紧随、复杂修饰后置”的尾重原则（End-Weight），通过 36 条核心语法法则（G001--G036）和 24 组高分功能句型模板（EP001--EP024）重塑句法骨架。
  \item \textbf{篇章信息流与图尔敏微论证体系（Given-to-New Flow \& Toulmin Argumentation）}：严密实施“已知信息在前、新信息在后”的主位推进流（Given-to-New Theme Progression）与 Toulmin 论证模型（Claim, Data, Warrant, Backing, Rebuttal, Qualifier），系统杜绝机械连接词与 AI 空洞套话（L001--L030）。
  \item \textbf{深度病段诊所与母语级重构图谱（Diagnostic Revision Clinic \& Full-Text Anthology）}：提供 60 例涵盖真实学科场景的典型病段，通过七维病灶透视与四步重构手术实现 publication-ready 升华；精选 24 篇跨五大学科门类的完整学术范文（FT01--FT24），以段落蓝图链（BP01--BP48）深度解构全篇气韵。
\end{enumerate}

\section*{How to Use This Book / 本书使用指南与学习路线图}
本书采用模块化、双轨制编排，既可作为高校研究生与科研人员的学术写作系统教材，亦可作为高强度日常打卡训练的实战练习册。

\begin{ConceptBox}{Visual Hierarchy \& Callout Legend / 排版视觉规范说明}
\begin{tabularx}{\linewidth}{lX}
  \textbf{\color{AcadNavy}ConceptBox (深海蓝)} & 核心理论定义、句法骨架法则、篇章论证机制与模型架构总览。\\
  \textbf{\color{AcadTeal}ExampleBox (雅致青)} & 真实高水平论文例句、经典发表级句式、病段与定稿对照呈现。\\
  \textbf{\color{AcadAmber}TipBox (暖琥珀)} & 常见学术陷阱警示、修辞去套话技巧、句法成分深度剖析。\\
  \textbf{\color{AcadForest}ExerciseBox (常青绿)} & 课后实战演练、改写训练、每日打卡任务。\\
  \textbf{\color{TargetHighlight}\textbackslash DefectTarget} & 病段中的核心语病处高亮标记（玫瑰红粗体）。\\
  \textbf{\textbackslash writinglines[11mm]} & 为手写习惯专门预留的 11mm 宽行距水平书写横线。\\
\end{tabularx}
\end{ConceptBox}

\section*{Four Specialized Study Tracks / 四大学习路径建议}
\begin{itemize}[leftmargin=1.5em]
  \item \textbf{路径 A：硕博论文发表快速攻坚（Fast-Track Publication）}：重点精读 Part II（句法架构）、Part III（篇章逻辑）、Part IV（48 段落蓝图）与 Part V（60 病段诊所），快速掌握学术正文各部分的标准段落展开套路与自查改写手术。
  \item \textbf{路径 B：60 天全景手写实战训练（Systematic Daily Practice）}：按 Part VII 练习册规划，每日完成填空、句子改写（SentenceRewriteItem）及段落重构（ParagraphRevision），在 11mm 书写线上扎实锤炼手写肌肉记忆，并每日对照 Part VIII 书末答案核验解析。
  \item \textbf{路径 C：国际会议演讲与答辩申请（Presentation \& Rebuttal）}：精读 Part IV Track 6（蓝图 BP41--BP48）及 Part VI 第 19--20 章（Oral 演讲稿 FT07--FT08、ICLR 审稿回复 FT09、答辩应答全案 FT10、个人陈述 SoP FT11、课题申请书 FT12）。
  \item \textbf{路径 D：语言考试高分冲刺（IELTS 8.0+ / TOEFL 110+ / 六级 600+）}：研读 Part VI 第 21--22 章中的雅思小作文（FT13--FT15）、托福综合写作（FT16）、六级写作（FT19--FT20）及雅思/托福大作文（FT21--FT24）。
\end{itemize}
\clearpage